\paragraph{Theorem E.2.6 (Proved: full integral homology of the shifted product).}
For the transverse shifted product $P$ of E.2.3, write
\[
 A=(\mathbb Z/12)^7,\qquad
 B=(\mathbb Z/12)^{10}\oplus(\mathbb Z/2)^7,\qquad
 C=(\mathbb Z/12)^6\oplus\mathbb Z/24\oplus(\mathbb Z/2)^{17}.
\]
Its complete integral homology is
\[
 \begin{array}{c|l}
 i&H_i(P;\mathbb Z)\\\hline
 0,1,8,9,10&\mathbb Z\\
 2&\mathbb Z\oplus A\\
 3&\mathbb Z^{20}\oplus B\\
 4&\mathbb Z^{63}\oplus C\\
 5&\mathbb Z^{88}\oplus C\\
 6&\mathbb Z^{63}\oplus B\\
 7&\mathbb Z^{20}\oplus A.
 \end{array}
\]
In particular the higher integral filtration extensions are nontrivial:
the middle torsion includes a cyclic factor of order $24$.
\begin{proof}
\emph{A decomposition with free groups on every piece.}
Choose a separating circle in $B=\mathbb P^1$ so that one closed
disc contains the three critical values of $X$ and the other
contains the three critical values of $X'$. The smooth torus
family of the other factor is trivial over each disc. Set
\[
 F=T^4,\qquad E=X\setminus\operatorname{int}(F\times D^2),
 \qquad E'=X'\setminus\operatorname{int}(F'\times D^2).
\]
The product trivializations give the actual decomposition
\[
 P=(E\times F')\ \cup_{F\times F'\times S^1}\ (F\times E').
 \tag{E.MV1}
\]
The boundary identification preserves the two marked torus factors
and reverses the normal circle, because the two base discs have
opposite boundary orientations. There is no additional torus
clutching: the same regular-disc trivializations define both the
boundary identifications of $E,E'$ and (E.MV1).

Since $X,X'$ are six-spheres, the Mayer--Vietoris sequence of
$S^6=E\cup(F\times D^2)$ gives, integrally,
\[
 H_0(E)=\mathbb Z,\qquad H_k(E)=H_{k-1}(F)\quad(1\leq k\leq4),
 \qquad H_k(E)=0\quad(k\geq5).
 \tag{E.MV2}
\]
The isomorphism is the normal-circle sweep
$\alpha\mapsto[S^1]\times\alpha$. Indeed, in degrees one
through four the boundary map to $H_k(E)\oplus H_k(F)$ is
an isomorphism, and its second component is the identity on
the unswept fibre summand and zero on the swept summand.
In degree five the boundary fundamental class is the image
of the sphere fundamental class, so $H_5(E)=0$.
Thus all homology groups of the two pieces and their common
boundary in (E.MV1) are free abelian. The same statements hold
for $E'$.

\emph{Determining every boundary map integrally.}
Let $V=H^1(F;\mathbb Z)$, in the basis $(\gamma,u,w,\delta)$,
and use the positively ordered volume
$v=\gamma u w\delta$. The primitive global invariant classes
of the original torus fibration are
\[
 I_0=1,\quad I_1=\gamma,\quad
 I_2=q=uw+6\gamma\delta,\quad
 I_3=\gamma uw,\quad I_4=v.
\]
The images of $H^k(E)\to H^k(F)$ are exactly
\[
 n_k\mathbb Z I_k,\qquad
 (n_1,n_2,n_3,n_4)=(12,2,2,12).
 \tag{E.MV3}
\]
Here is an integral verification, rather than descent of invariant
real forms. On the disc the only local restrictions on a global
invariant class are the actual finite-cover pullback lattices;
the cusp specialization is an integral isomorphism onto its
invariant lattice by E.2.5. In degree one the finite images
contain $3\gamma$ and $4\gamma$, giving the intersection
$12\mathbb Z\gamma$. In degree two the order-three image
contains $q$, while the order-four image requires an even
coefficient of $q$, by the finite integral lattices of B.5.13.
In degree three use the perfect integral pairing downstairs:
an invariant class is a pullback precisely when its pairings
with every pulled-back degree-one class are divisible by the
covering degree. The other degree-one generators are
$2u+w+3\delta$ and $u+w+2\delta$. Their pairings with
$t\gamma uw$ are, up to signs, $3t$ and $2t$.
Thus the first cover imposes no further condition, and the
second requires $t$ even. In degree four the pullback degrees
are three and four, giving $12\mathbb Z v$.

These local restrictions are also sufficient for the image
from $E$. For a local system on a disc with punctures, the
restriction from its free-group $H^1$ to the direct sum of
the local coinvariants is onto: each puncture has its own
independent free generator. Consequently
$H^2(D,j_*\mathcal L)=0$. Replacing $j_*\mathcal L$ by the
actual direct-image sheaf changes it only by point-supported
cokernels, and so does not change this vanishing.
The Leray sequence over the disc therefore makes
$H^k(E)\to H^0(D,R^kf_*\mathbb Z)$ onto. This proves
(E.MV3) without an integral formality assumption.

Write $\theta$ for the normal-circle covector. The dual of
(E.MV2) identifies $H^k(E)$ with $\bigwedge^{k-1}V$, and
its boundary restriction has the graph form
\[
 r_k(\alpha)=\phi_k(\alpha)+\theta\alpha,
 \qquad
 \phi_k:\bigwedge^{k-1}V\longrightarrow\bigwedge^kV.
 \tag{E.MV4}
\]
The image of $\phi_k$ is exactly $n_k\mathbb Z I_k$ by
(E.MV3). If $k+l=5$, classes extending to $E$ have zero
product evaluated on $\partial E$. Therefore
\[
 (-1)^k\phi_k(\alpha)\beta+
                    \alpha\phi_l(\beta)=0
 \quad\text{in }\bigwedge^4V.
 \tag{E.MV5}
\]
Each $I_k$ is primitive, and wedge pairing between the two
complementary exterior powers is unimodular. Writing
$\phi_k=f_k I_k$, equation (E.MV5), tested against a class
pairing to one with $I_k$, proves that $f_k$ is an integral
multiple of $\alpha\mapsto\langle\alpha I_{5-k},[F]\rangle$.
The exact image (E.MV3) fixes the absolute multiple. Thus
\[
 \begin{split}
 \phi_k(\alpha)&=c_k
       \langle\alpha I_{5-k},[F]\rangle I_k,\\
 (c_1,c_2,c_3,c_4)&=(12\varepsilon,2\eta,-2\eta,-12\varepsilon),
 \qquad \varepsilon,\eta\in\{1,-1\}.
 \end{split}
 \tag{E.MV6}
\]
This proves every boundary map up to two signs. We allow those
signs independently for $E$ and $E'$ below; the answer is the
same for all sixteen possibilities.

\emph{The full integral Mayer--Vietoris matrices.}
Let $i_p:H_p(F)\to H_p(E)$ be the transpose of $\phi_p$ in
the sweep bases, with $i_0=1$; let $j_{p+1}$ denote the sweep
isomorphism for $0\leq p\leq3$, and zero for $p=4$.
The analogous maps of the second copy are $i'_q,j'_{q+1}$.
For a fibre tensor $\alpha\otimes\beta$, with degrees $p,q$,
the matrix $D_k$ of (E.MV1) is
\[
 \alpha\otimes\beta\longmapsto
  (i_p\alpha\otimes\beta,-\alpha\otimes i'_q\beta),
 \qquad p+q=k.
 \tag{E.MV7}
\]
For a swept tensor, place the normal circle first. Then
\[
 \theta\times\alpha\times\beta\longmapsto
 \bigl(j_{p+1}\alpha\otimes\beta,
                 (-1)^p\alpha\otimes j'_{q+1}\beta\bigr),
 \qquad p+q=k-1.
 \tag{E.MV8}
\]
The second sign contains the Koszul sign from moving the circle
past $\alpha$, the reversal of the two boundary-circle
orientations, and the second minus in Mayer--Vietoris. The last
two cancel. This specifies every integer entry of $D_k$.

The executable \texttt{code/manifolds\_shifted\_mv.py} constructs
the exterior bases, the four adjoint maps (E.MV6), and the full
matrices (E.MV7)--(E.MV8). Exact Smith reduction gives the table
below. The source ranks are $\binom9k$; a zero binomial is
understood outside $0\leq k\leq9$.
\[
 \begin{array}{c|rrrr|l}
 k&\text{target rank}&\text{source rank}&\operatorname{rank}D_k&
       \text{unit factors}&\text{nonunit factors}\\\hline
 0&2&1&1&1&--\\
 1&10&9&9&9&--\\
 2&28&36&27&20&12^{[7]}\\
 3&64&84&53&36&2^{[7]},12^{[10]}\\
 4&114&126&82&58&2^{[17]},12^{[6]},24\\
 5&138&126&94&70&2^{[17]},12^{[6]},24\\
 6&104&84&73&56&2^{[7]},12^{[10]}\\
 7&44&36&35&28&12^{[7]}\\
 8&8&9&8&8&--\\
 9&0&1&0&0&--\\
 10&0&0&0&0&--
 \end{array}
\]
An exponent $[r]$ records repetition, not an integer power.
All sixteen choices of the two signs in each copy give precisely
this table; the retained output and JSON contain every case.

Finally the exact homology sequence gives
\[
 0\longrightarrow\operatorname{coker}D_k
 \longrightarrow H_k(P)
 \longrightarrow\ker D_{k-1}\longrightarrow0.
\]
The last group is free, since the boundary homology is free.
This sequence therefore splits as abelian groups, so the Smith
factors in the table are the actual torsion of $H_k(P)$.
Adding the free kernel ranks gives exactly the rational Betti
vector in E.2.3 and the integral groups asserted here.
The resulting torsion also satisfies the required ten-dimensional
duality $\operatorname{Tor}H_k(P)\cong
\operatorname{Tor}H_{9-k}(P)$.
\end{proof}

\paragraph{Proposition E.2.7 (Proved: a unique transverse family type and its stable tangent bundle).}
For fixed marked sphere fibrations, all M\"obius transformations
whose critical triples are disjoint give orientation-preservingly
diffeomorphic shifted products. The same conclusion holds while
either sphere varies in its proper smooth marked period family.
For every such fivefold $P$ there is an isomorphism of real vector
bundles
\[
 TP\oplus\mathbb R^4\simeq\mathbb R^{14}.
 \tag{E.MV9}
\]
Thus all its positive-degree Stiefel--Whitney and Pontryagin
classes vanish. For its given complex tangent bundle,
$c_1(P)=c_2(P)=0$ integrally.
These assertions and E.2.6 do not identify the complete smooth
structure from its ordinary homology and fundamental group.
\begin{proof}
The allowed transformations form the complement in
$\operatorname{PGL}_2(\mathbb C)$ of the nine analytic
hypersurfaces imposing an equality of one critical value from
each triple. This complement is path connected: a piecewise
smooth path between two points can be perturbed relative to
its endpoints to avoid these sets, which have real codimension
two. Equivalently, use local complex coordinate balls and
perturb each segment there. Choose a smooth such path.

The corresponding fibre products form a compact smooth family
over its parameter interval. Indeed the vertical derivative of
$f_1(x)-f_2(y)$ is surjective at every point, since its two
factors cannot be critical at the same base value. The same
surjectivity solves the equation for a lift of the interval
coordinate vector field. A partition of unity supplies a
global smooth lift; compactness makes its flow exist for the
whole interval. This flow gives the claimed diffeomorphisms.
The complex orientations vary continuously and hence are
preserved. The identical argument applies to smooth proper
period families of either factor, keeping the critical triples
disjoint throughout.

View $P$ as the transverse inverse image of the diagonal in
$B\times B$ under $f_1\times f_2$. Its real oriented normal
two-plane bundle is the underlying bundle of $g^*T_B$,
where $g:P\to B$ is the common base map. It is trivial:
its first Chern class is
\[
 c_1(g^*T_B)=2g^*[B]^\vee
    =2\operatorname{pr}_1^*f_1^*[B]^\vee=0,
\]
because $H^2(X;\mathbb Z)=0$, and complex line bundles are
classified by their first Chern class. The tangent bundles
of the two six-spheres become trivial after adding one real
line each. Restriction of their product tangent bundle gives
\[
 (TP\oplus\mathbb R^2)\oplus\mathbb R^2
       \simeq\mathbb R^{14},
\]
proving (E.MV9) and the characteristic-class assertions for
the real tangent bundle. Finally the holomorphic tangent
sequence is
\[
 0\longrightarrow TP\longrightarrow
 \operatorname{pr}_1^*TX\oplus\operatorname{pr}_2^*TX'
 \longrightarrow g^*T_B\longrightarrow0.
\]
Both sphere factors have integral $c_1=c_2=0$, since their
cohomology in degrees two and four vanishes. The last line
has zero first Chern class, as just proved. Multiplying total
Chern classes in this exact sequence gives
$c_1(P)=c_2(P)=0$ integrally.
\end{proof}

\paragraph{Proposition E.2.8 (Proved: the shifted product bounds with its natural stable framing).}
Equip the two oriented six-spheres with their standard bounding
stable tangent framings. Trivialize the oriented two-plane bundle
$f_1^*T\mathbb P^1$ on the first sphere, and equip the transverse
shifted product $P$ with the induced stable framing in (E.MV9).
Then there is a compact smooth oriented eleven-manifold $V$ with
\[
 \partial V=P,\qquad TV\oplus\mathbb R^3\simeq\mathbb R^{14},
 \qquad [P]=0\quad\hbox{in }\Omega^{\mathrm{fr}}_{10}.
 \tag{E.B1}
\]
The boundary framing is the stated one. In particular its induced
spin structure bounds. This supplies a framed bordism invariant;
it does not identify the smooth type from ordinary homology.
\begin{proof}
Write $f_2:S^6\to\mathbb P^1$ for the shifted map of the second
factor. By B.2.1 the first map $f_1:S^6\to\mathbb P^1$ is
null-homotopic. It therefore extends to a smooth map
$F:D^7\to\mathbb P^1$, which can be chosen to equal $f_1$ on
each sphere of a collar of the boundary. To see this with the
boundary fixed, first use the null-homotopy to obtain a continuous
extension constant in the collar direction. Smooth its interior
in coordinate charts, with cutoffs supported away from the collar;
on the overlap the extension is already smooth.

Let $C_2$ be the three critical values of $f_2$. The restriction
$f_1$ is transverse to each point of $C_2$, since the two critical
triples are disjoint. Perturb $F$ only in the interior to be
transverse to all these points. The applicable finite simultaneous
transversality statement is
\cite[Chapter 3, Section 2, Theorems 2.1 and 2.5]{rootHirsch76}.
For the relative assertion, the proof of that theorem uses local
perturbations multiplied by cutoff functions. Their supports may
be restricted to the complement of the collar, where
transversality has not already been secured; the open condition
on the collar persists. Compactness permits finitely many such
perturbations. Thus no assumption that $F$ is a submersion
everywhere is needed.

Define the smooth fibre product
\[
 V=\{(x,y)\in D^7\times S^6:F(x)=f_2(y)\}.
 \tag{E.B2}
\]
At a point whose common value is outside $C_2$, the differential
$df_2$ is onto. At a point over $C_2$, the differential $dF$ is
onto. Consequently $dF-df_2$ is onto everywhere on (E.B2).
The same argument holds on its boundary, using $df_1$ there.
The implicit-function theorem, including its boundary version,
makes $V$ a smooth eleven-manifold with boundary exactly $P$.
It is compact as a closed subset of $D^7\times S^6$.
The inverse-image orientation using the complex orientation of
$T\mathbb P^1$ induces the complex orientation of $P$ at the
boundary, because the collar keeps the first map fixed.

Let $g:V\to\mathbb P^1$ be the common map. The real normal
two-plane bundle of (E.B2) is $g^*T\mathbb P^1$, identified by
$dF-df_2$. This is the pullback from $D^7$ of $F^*T\mathbb P^1$,
which is trivial because the disc is contractible. A choice of
its oriented framing restricts on the boundary to the chosen
framing of $f_1^*T\mathbb P^1$ up to homotopy: two oriented
framings of this two-plane bundle differ by a map
$S^6\to SO(2)$, and $H^1(S^6;\mathbb Z)=0$.
The tangent exact sequence and the bounding stable framing of
the second sphere now give
\[
 \begin{split}
 TV\oplus\mathbb R^2&\simeq
       T(D^7\times S^6)|_V,\\
 TV\oplus\mathbb R^3&\simeq
       \mathbb R^7\oplus(TS^6\oplus\mathbb R)|_V
       \simeq\mathbb R^{14}.
 \end{split}
\]
At the boundary, the outward normal of $D^7$ supplies the
additional real line in $TV|_P=TP\oplus\mathbb R$.
The resulting framing is precisely the stable tangent framing
obtained from the two sphere framings and the normal framing
in E.2.7. This proves (E.B1) by the definition of framed bordism.
It also proves the assertion about the induced spin structure,
by forgetting the framing.

This disc-extension construction does not by itself extend a
classifying map $P\to B\pi_1(P)$ over $V$.
The different filling in E.2.9 below does extend that map.
Both are bordism statements; the higher equivariant surgery
data required for the smooth classification remain separate.
\end{proof}


\paragraph{Theorem E.2.9 (Proved: framed null-bordism over the fundamental group).}
Let $G$ be the group in (E.16). There is a closed smooth
aspherical four-manifold $N$ with
\[
 \pi_1(N)=G,\qquad \widetilde N\simeq\mathbb H^2\times\mathbb R^2,
 \qquad H_*(N;\mathbb Z)\cong H_*(S^1\times S^3;\mathbb Z),
 \tag{E.B3}
\]
and a smooth classifying map $c:P\to N$ inducing the displayed
identification of fundamental groups. With the natural stable
framing of E.2.8, one has
\[
                 [P,c]=0\quad\hbox{in }\Omega^{\mathrm{fr}}_{10}(N).
 \tag{E.B4}
\]
Thus $P$ has a compact stably framed eleven-dimensional filling
over $BG$. This is stronger than the filling of E.2.8; it is
still a bordism statement, not a diffeomorphism classification.
\begin{proof}
\emph{The aspherical target.}
The repaired quotient maps of A.2.1 are smooth maps
$e_i:S^6\to S^3$ over the base, and B.2.1 gives the exact
factorizations $f_i=h_i e_i$. Here $h_i:S^3\to\mathbb P^1$
is the weighted Hopf map of orders $3,4$, with the second
map shifted by the chosen M\"obius transformation. Define
\[
 N=\{(u,v)\in S^3\times S^3:h_1(u)=h_2(v)\},\qquad
 c(x,y)=(e_1(x),e_2(y)).
 \tag{E.B5}
\]
The critical values of each $h_i$ are only its two exceptional
circle values. They are disjoint, so the same transversality
argument as in E.2.2 makes $N$ a compact smooth four-manifold.
The two circle actions on the weighted Hopf spheres give a
principal $T^2$ bundle in the orbifold sense over
\[
                  \mathcal B=S^2(3,4,3,4).
\]
At a cone of order $m$, its uniformizing chart is
$(D^2\times T^2)/C_m$, with the generator rotating the disc
and translating one torus coordinate by a primitive fraction
of denominator $m$. This local description is smooth and is
the same finite-filling description as A.1.3, with the unused
circle retained.

The orbifold has Euler characteristic
$2-2(1-1/3)-2(1-1/4)=-5/6$. It is good and has hyperbolic
orbifold universal cover, by
\cite[Theorems 2.3--2.4 and the Euler-number discussion in
Section 2]{rootScott83}. Pulling the principal torus bundle to
$\mathbb H^2$ gives an ordinary principal $T^2$ bundle over a
contractible base, hence the trivial bundle. Its universal
cover is $\mathbb H^2\times\mathbb R^2$. The local chart above
shows that these orbifold pullbacks are ordinary covering maps
of total spaces. This proves asphericity and injectivity of
the fibre subgroup $\mathbb Z^2$.

Using the marked fibre generators $h_1,h_2$ and the four
positive peripheral generators gives exactly
\[
 \begin{split}
 \pi_1(N)=\langle a_1,a_2,b_1,b_2,h_1,h_2\mid {}&h_1,h_2
       \text{ central},\quad a_1^3=h_1,\quad a_2^4=h_1^{-1},\\
 &b_1^3=h_2,\quad b_2^4=h_2^{-1},\quad
 a_1a_2b_1b_2=1\rangle.
 \end{split}
\]
The map $c$ sends the surviving regular-fibre generators
$\widehat\gamma,\widehat\gamma'$ to $h_1,h_2$, respectively,
and each peripheral generator to the one with the same name.
This is because $e_i$ is the literal circle quotient on these
fibres and finite collars. Its cusp replacement extends over
a solid torus and kills precisely the cusp meridian in this
presentation. Comparing with the full van Kampen calculation
in E.2.3 therefore shows that $c_*$ is the identity on the
displayed presentation, and in particular is an isomorphism.
As $N$ is aspherical, $c$ is a classifying map.

For completeness the abelian relation matrix has Smith
diagonal $(1,1,1,1,1)$; a primitive integer kernel vector in
the order $(a_1,a_2,b_1,b_2,h_1,h_2)$ is
$(4,-3,-4,3,12,-12)$. Thus $H_1(N)=\mathbb Z$.
Every stratum of its base has a torus fibre, so Euler
additivity gives $\chi(N)=0$. Integral Poincare duality and
the universal coefficient theorem give $H_3(N)=\mathbb Z$
and no torsion in $H_2(N)$. Euler characteristic then forces
$H_2(N)=0$, proving the integral homology assertion in (E.B3).

\emph{The bordism computation before taking the inverse image.}
For a based sphere, the stable splitting
$(S^3)_+\simeq S^0\vee S^3$ identifies
\[
 \Omega^{\mathrm{fr}}_6(S^3)
       =\pi^{\mathrm S}_6\oplus\pi^{\mathrm S}_3.
 \tag{E.B6}
\]
We use Pontryagin--Thom in its form for maps to a space,
\cite[Definition 8.16 and Theorem 8.17]{rootDavisKirk01}.
The first component of $[S^6,e_i]$ is obtained by forgetting
the map, so it is zero for the bounding sphere framing.
The second component is the stable class of $e_i$, equivalently
the framed inverse image of a regular point. It is $\eta^3$
by B.1.2. One can also check the splitting directly: the
Pontryagin--Thom collapse of the framed sphere into its
disjoint-basepoint suspension has degree one on the sphere
summand; composing with $e_i$ gives its stable homotopy class.
Thus
\[
                         [S^6,e_i]=(0,\eta^3).
\]
Smashing the two sphere splittings gives
\[
 (S^3\times S^3)_+\simeq
       S^0\vee S^3\vee S^3\vee S^6.
\]
Under this splitting the external product
$[S^6\times S^6,e_1\times e_2]$ has four components.
The first three each contain a factor $[S^6]=0$; the last
is $\eta^3\eta^3=\eta^6=0$. Indeed $\eta^4$ belongs to
$\pi^{\mathrm S}_4=0$, as recorded in
\cite[Theorem 2.1 and Table 1]{rootIWX20}. Therefore
\[
 [S^6\times S^6,e_1\times e_2]=0
      \quad\hbox{in }\Omega^{\mathrm{fr}}_{12}(S^3\times S^3).
 \tag{E.B7}
\]
This calculation uses the product of the specific bounding
sphere framings, not just vanishing of ordinary characteristic
classes.

\emph{The comparison with the requested classifying map.}
By (E.B7) there is a compact stably framed thirteen-manifold
$Z$ with boundary $S^6\times S^6$ and a map
$E:Z\to S^3\times S^3$ extending $e_1\times e_2$.
Smooth it relative to the boundary and make it transverse
to $N$ there. Transversality on the boundary already holds:
the derivative of its defining equation is $df_1-df_2$,
which is onto. Relative transversality therefore supplies
the interior perturbation without changing that boundary.

The normal bundle of $N\subset S^3\times S^3$ is the
restriction of $h_1^*T\mathbb P^1$ pulled back from the first
sphere. This oriented two-plane bundle is trivial because
$H^2(S^3;\mathbb Z)=0$. Choose a framing on it before
restriction to $N$. Its pullback under $e_1\times e_2$
is a framing of $f_1^*T\mathbb P^1$, agreeing with that
in E.2.8 up to homotopy since $H^1(S^6;\mathbb Z)=0$.
It follows that
\[
                  V=E^{-1}(N),\qquad E|_V:V\longrightarrow N
\]
is a compact stably framed eleven-manifold with boundary
$(P,c)$ and the required boundary framing. This proves (E.B4).

The exact group-matrix and Euler calculation, together with
the two equivalent stable-product reductions, is recorded in
\texttt{code/shifted\_group\_bordism.py}.
The earlier disc-extension filling need not extend the
classifying map; (E.B7) supplies a different filling that does.
Neither construction computes the higher equivariant surgery
data needed to distinguish smooth ten-manifolds over $N$.
\end{proof}


\input{sections/maps_shifted_homotopy}
